\section{Rancangan Pipeline Konstruksi Basis Pengetahuan}

\subsection{Pendekatan Modular}

Paper ini memandang \plugin sebagai unit penggantian strategi, bukan sebagai topik implementasi tersendiri. Orkestrator hanya mengatur urutan tahap dan membaca artefak dengan tipe yang sama; keputusan tentang cara membentuk chunk, indeks, atau graf ditempatkan pada komponen yang dipilih. Dengan batas tersebut, perbandingan antarkomponen dapat dilakukan tanpa mengubah tahap hilir secara bersamaan.

Satu dokumen hasil parser menjadi sumber bersama untuk tiga proyeksi: unit pencarian melalui \chunker, representasi pencarian melalui \indexer, dan hubungan antardokumen melalui \kg. Representasi kanonik menjadi titik temu ketiganya. Provenance dan metadata dipertahankan sebagai bagian dari artefak sehingga hasil tiap strategi dapat ditelusuri kembali ke dokumen serta konfigurasi asalnya.

\begin{table}[t]
\centering
\caption{Peran artefak pada pendekatan modular}
\label{tab:kontrak-artefak}
\scriptsize
\begin{tabularx}{\columnwidth}{@{}p{0.31\columnwidth}X@{}}
\toprule
\textbf{Artefak} & \textbf{Peran dalam perbandingan}\\
\midrule
\code{CanonicalDocument} & Dokumen terurai yang memuat teks, hierarki, lokasi sumber, dan metadata; menjadi masukan semua strategi \chunker.\\
\code{CanonicalChunks} & Unit retrieval dan context dengan identitas posisi serta provenance; dipakai bersama oleh \indexer dan \kg.\\
\code{CanonicalIndex} & Proyeksi lexical, dense, atau hybrid dari chunk; menjadi masukan jalur direct retrieval.\\
\code{CanonicalGraph} & Proyeksi simpul dan relasi beserta bukti sumber; dipakai opsional untuk memperluas lingkungan dokumen.\\
\bottomrule
\end{tabularx}
\end{table}

Kontrak antarmuka menjaga agar keluaran setiap strategi dapat dikonsumsi tahap berikutnya dengan bentuk yang konsisten. Pemilihan strategi dilakukan melalui konfigurasi pipeline, sedangkan provenance mencatat komponen, versi, dan artefak yang digunakan. Dengan demikian, manfaat \plugin yang diuji dalam paper ini adalah keterpisahan faktor eksperimen: perubahan satu strategi tidak perlu disertai perubahan pada orkestrator atau format data downstream.

\subsection{Komponen yang Dibandingkan}

Kelompok \chunker membandingkan \textit{sliding-window}, \textit{recursive}, \textit{generic structure-aware}, dan \textit{legal structure-aware}. \textit{Sliding-window} menjadi baseline berbasis jendela tetap; \textit{recursive} memanfaatkan pemisah teks bertingkat; strategi structure-aware membawa heading atau hierarki hukum sebagai konteks. Pada dokumen hukum, parent chunk menjaga unit Pasal yang utuh sedangkan child chunk menjadi kandidat pencarian ketika unit tersebut terlalu panjang.

Kelompok \indexer membandingkan representasi sparse, dense, dan hybrid dari \code{CanonicalChunks}. Sparse mengandalkan kecocokan lexical, dense menggunakan embedding BAAI/BGE-M3, dan hybrid menggabungkan dua daftar kandidat dengan Reciprocal Rank Fusion. Perbedaan yang dibaca pada hasil adalah perbedaan representasi dan urutan kandidat, bukan perbedaan dokumen masukan.

Kelompok \kg menggunakan legal graph deterministik sebagai proyeksi opsional. Graf memuat hubungan struktur dan perubahan peraturan, lalu menerima seed dari hasil top-$K$ retrieval. Traversal dibatasi oleh kedalaman, arah, jenis relasi, dan cakupan eksperimen; hasilnya dipetakan kembali ke chunk atau dokumen yang memiliki provenance. Karena itu, graf memperluas lingkungan dokumen tanpa menggantikan text retrieval sebagai sumber bukti utama.

\subsection{Cara Pendekatan Diuji}

Setiap kelompok membandingkan satu faktor pada satu waktu. Ketika \chunker diuji, datasource, parser, embedding, dan \indexer dikunci. Ketika \indexer diuji, \chunker dan seluruh input dikunci. Ketika \kg diuji, seed retrieval serta aturan traversal ditetapkan sebelum pengukuran. Direct retrieval mematikan reranker dan graf agar metrik retrieval tidak tercampur dengan tahap hilir.

Tiga klaim operasional kemudian diuji: (1) pemeliharaan struktur meningkatkan cakupan context reference dibandingkan fixed-window; (2) dense retrieval meningkatkan recall ketika kueri memparafrasekan istilah sumber, sementara hybrid dapat menguntungkan posisi hit pertama; dan (3) traversal legal graph memperluas jangkauan dokumen tetangga, tetapi neighborhood recall harus dibaca terpisah dari context-reference recall. Hasil tiap konfigurasi dilaporkan pada cutoff dan korpus yang sama sehingga perbedaan dapat dikaitkan dengan strategi yang dibandingkan.
